\section{The common path numerator and positive arithmetic classes}
\label{sec:v74-normalization}
The scalar source estimate need not be repeated for each bounded
path mark. We give the precise downstream implication, keeping
its analytic input visible.

Let $\mathcal X=C([0,1],\R^4)$, and use the bounded-Lipschitz class
with supremum norm and Lipschitz constant at most one, including
the constant function one. The notation $\|\cdot\|_{\mathrm{BL}^*}$
denotes its dual norm. The common roof kernel and a countable
norming class fix all versions before taking essential suprema.
Write
\[
 \mathbf P(t)=m^2p(t)\mathsf Q_t,\qquad P(t)=\mathbf P(t)(1),
 \qquad G(t)=\mathcal L_{m,R}(k_1,k_2,t,n).
\]
The reference path probability $\mathsf W_R$ is the one in the
inherited path theorem, with its original covariance and arithmetic
coefficient. For each fixed $B$ the assumed, source-pinned theorem
states
\[
 \mathbf P-G\mathsf W_R=\mathbf b_B+\mathbf E_B,
 \qquad \mathbf b_B\ge0,\quad
 \mathbf b_B(1)=m^2b^{\varepsilon(B)}.
\]
The ordered error norm of $\mathbf E_B$ is at most $C_MB^{-1/192}$.
This is the inherited positive path-remainder theorem, including
its cylinder inversion and tightness inputs; it is not inferred
from the scalar $BV$ theorem.

\begin{theorem}[Equality and normalization on the same roofs]
\label{thm:v74-path}
Under this positive path-remainder input,
$\mathcal E_M^{\mathrm{path}}=\mathcal E_M$, in the extended
nonnegative reals. If the scalar endpoint vanishes, then on every
original central target with $G\ge d>0$,
\[
 d_{\mathrm{BL}}(\mathsf Q_t,\mathsf W_R)\longrightarrow0,
 \qquad P/G\longrightarrow1,\qquad G/P\longrightarrow1
\]
uniformly in essential supremum, radius and exact labels within
that target. On any positive-length roof window in the same target,
the forward normalized arithmetic likelihood tends to one in
essential supremum. No normalization is made on zero arithmetic
classes. The path assertion is in bounded-Lipschitz dual norm,
not total variation on path space.
\end{theorem}
\begin{proof}
Put $e_B=\|\mathbf E_B\|_{\mathrm{BL}^*}$. Positivity and the
constant test give $\|\mathbf b_B\|_{\mathrm{BL}^*}=\mathbf b_B(1)$,
and hence
\[
 |P-G|\le\|\mathbf P-G\mathsf W_R\|_{\mathrm{BL}^*}
                                      \le |P-G|+2e_B.
\]
Take the collision limsup at fixed $B$, then let $B\to\infty$.
This proves the equality, including the infinite case by the
left inequality.

Let $a=|P-G|$ and $D=\|\mathbf P-G\mathsf W_R\|_{\mathrm{BL}^*}$.
When $G\ge d$ and $a\le d/2$, one has $P\ge d/2$ and
\[
 d_{\mathrm{BL}}(\mathsf Q_t,\mathsf W_R)\le\frac{D+a}{P}
       \le\frac{2(D+a)}d,\qquad
 |P/G-1|\le a/d,\qquad |G/P-1|\le2a/d.
\]
These inequalities imply the stated limits. For a roof window $J$
of positive length put $F_J=\int_J P$, $H_J=\int_J G$ and
$r=\|P-G\|_{L^\infty(J)}/d$. Then
$|F_J/H_J-1|\le r$. For $r<1$ the normalized original roof laws
satisfy
\[
 \left\|\frac{d\mathbb P_J}{d\mathbb Q_J}-1\right\|_\infty
   =\left\|\frac{P/G}{F_J/H_J}-1\right\|_\infty
                                      \le\frac{2r}{1-r}.
\]
All ratios are taken only where the reference floor holds. No
absolute continuity between path-space probabilities is asserted.
\end{proof}

The actual-return path additionally uses the inherited
collision-to-return clock theorem on these same positive classes.
The preceding argument supplies its collision numerator after the
scalar endpoint; it introduces no new clock. Thus a proof of the
full-source short-excursion decay in
\eqref{eq:v74-error-limit}, for every central $M$, would propagate
to the original downstream conclusions. The decay and an independent
specialist audit of the inherited spectral and clock inputs are not
supplied by the new deterministic inequalities.
